\documentclass[11pt]{article}

\usepackage[margin=1in]{geometry}
\usepackage{amsmath, amssymb}
\usepackage{xcolor}
\usepackage[skins,breakable]{tcolorbox}
\usepackage{hyperref}

\hypersetup{
    colorlinks=true,
    linkcolor=blue!60!black,
    urlcolor=blue!60!black,
    citecolor=blue!60!black
}

\newcommand{\SymPy}{SymPy}

% Public artifact repository; \bugbundle links a bug's bundle folder into it.
% TODO: set the public artifact repository URL before release.
\newcommand{\artifactrepo}{https://github.com/TODO-org/TODO-repo}
\newcommand{\bugbundle}[1]{\href{\artifactrepo/tree/main/bug-report/#1}{\nolinkurl{bug-report/#1/}}}

\newtcolorbox{counterexamplebox}{
  enhanced,
  colback=yellow!5,
  colframe=orange!70!black,
  arc=2mm,
  boxrule=0.8pt,
  left=4mm,
  right=4mm,
  top=3mm,
  bottom=3mm,
  before=\par\medskip\noindent,
  after=\par\medskip,
  before upper=\raggedright
}

\title{SymPy Bug Card}
\author{}
\date{}

\begin{document}

\maketitle

\section*{Bug 42: \texorpdfstring{\texttt{solve\_univariate\_inequality(atan(x) < pi)}}{solve\_univariate\_inequality(atan(x) < pi)} omits \texorpdfstring{\(x=0\)}{x=0}}

\begin{counterexamplebox}
\textbf{Task.} Solve the real univariate inequality
\[
\arctan(x)<\pi .
\]

\medskip
\textbf{SymPy's answer.}
\[
\operatorname{solve\_univariate\_inequality}(\arctan(x)<\pi,x)
  =(-\infty,0)\cup(0,\infty).
\]

\medskip
\textbf{Correct answer.}
\[
\mathbb{R}.
\]

\medskip
\textbf{Explanation.} The returned set excludes \(0\), but
\(\arctan(0)=0<\pi\).  More generally, real principal arctangent values lie in
\((-\pi/2,\pi/2)\), so the inequality is true for every real \(x\).

\medskip
\textbf{Likely root cause.} The diagnosis identifies a branch-invalid inversion
in \texttt{solveset}: \(\tan(\pi)=0\) is accepted as a root of
\(\arctan(x)=\pi\), and the inequality solver then removes this spurious root as
a strict boundary.  This is a new failure mode with no public precedent.

\medskip
\textbf{Artifact (GitHub).} \bugbundle{bug-042-solve-univariate-inequality-atan-x-pi-omits-x-0}
\end{counterexamplebox}

\end{document}
